\documentclass[11pt]{article}

\usepackage{epsfig}
\usepackage{amsfonts}
\usepackage{amssymb}
\usepackage{amstext}
\usepackage{amsmath}
\usepackage{xspace}
\usepackage{theorem}
\usepackage{xcolor}
\usepackage{graphicx}
\usepackage{url}


% This is the stuff for normal spacing
\makeatletter
 \setlength{\textwidth}{6.5in}
 \setlength{\oddsidemargin}{0in}
 \setlength{\evensidemargin}{0in}
 \setlength{\topmargin}{0.25in}
 \setlength{\textheight}{8.25in}
 \setlength{\headheight}{0pt}
 \setlength{\headsep}{0pt}
 \setlength{\marginparwidth}{59pt}

 \setlength{\parindent}{0pt}
 \setlength{\parskip}{5pt plus 1pt}
 \setlength{\theorempreskipamount}{5pt plus 1pt}
 \setlength{\theorempostskipamount}{0pt}
 \setlength{\abovedisplayskip}{8pt plus 3pt minus 6pt}
\makeatother


%%%%%%%%%%%%%%%%%%%%%%%%%%%%%%%%%%%%%%%%%%%%%%%%%%%%%%%%%%%%%%%%%%%%%%%%%%%
% Document begins here %%%%%%%%%%%%%%%%%%%%%%%%%%%%%%%%%%%%%%%%%%%%%%%%%%%%
%%%%%%%%%%%%%%%%%%%%%%%%%%%%%%%%%%%%%%%%%%%%%%%%%%%%%%%%%%%%%%%%%%%%%%%%%%%

\newcommand{\headings}[3]{{\large
{\bf CSE548: Analysis of Algorithms, Fall 2026} \hfill {{\bf Release Date:} #1}\\
\rule[0.01in]{\textwidth}{0.025in}}
\begin{center} {{\bf\huge{Homework \##2}}\\{\bf( Due Date: #3 )}} \end{center}
}

\begin{document}

\headings{Oct 25}{2}{Nov 6}
%\headings{{\underline{\hspace{1.5cm}\textcolor{white}{.}}}}{{\underline{\hspace{1cm}}}}{{\underline{\hspace{1.6cm}}}}
\newcommand{\lecnum}{5}  % Lecture Number


\vspace{2in}

{\Large
\begin{center}

{\sc{Group Number:} {\underline{17}}}~\\~\\

\scalebox{0.83}{
\begin{tabular}{||l|c|c||}\hline\hline
 \multicolumn{3}{||c||}{Group Members}\\ \hline
Name\hspace{3in}                              &      \hspace{0.5in}SBU ID\hspace{0.5in}     &     \% Contribution \\ \hline\hline
        Lichen Yang      &       116472922          &       100 \%         \\
                                  &                 &                     \\ \hline
                                  &                 &                     \\
                                  &                 &                     \\ \hline
                                  &                 &                     \\
                                  &                 &                     \\ \hline\hline
\end{tabular}
}
\end{center}
}


\newpage

{\Large
\begin{center}

{\sc{\underline{Collaborating Groups}}}~\\
\vspace{0.25cm}
\scalebox{0.83}{
\begin{tabular}{||c|c||}\hline\hline
Group         &  Specifics \\
Number        &  (e.g., specific group member? specific task/subtask?) \\ \hline\hline
    N/A     &                    None                  \\ 
              &                                      \\ 
              &                                      \\ \hline
              &                                      \\ 
              &                                      \\ 
              &                                      \\ \hline
              &                                      \\ 
              &                                      \\ 
              &                                      \\ \hline
              &                                      \\ 
              &                                      \\ 
              &                                      \\ \hline
              &                                      \\ 
              &                                      \\ 
              &                                      \\ \hline\hline
\end{tabular}
}
\end{center}
}



{\Large
\begin{center}

{\sc{\underline{External Resources Used}}}~\\
\vspace{0.25cm}
\scalebox{0.83}{
\begin{tabular}{||c|c||}\hline\hline
              &  Specifics \\
              &   (e.g., cite papers, webpages, etc. and mention why each was used) \\ \hline\hline
              & OpenAI Codex: writing, polishing, and verification. \\  
    $1.$      &                                      \\ 
              &                                      \\ \hline
              &                                      \\ 
    $2.$      &                                      \\ 
              &                                      \\ \hline
              &                                      \\ 
    $3.$      &                                      \\ 
              &                                      \\ \hline
              &                                      \\ 
    $4.$      &                                      \\ 
              &                                      \\ \hline
              &                                      \\ 
    $5.$      &                                      \\ 
              &                                      \\ \hline\hline
\end{tabular}
}
\end{center}
}



\newpage
\section*{Task 1. A Deterministic Sorting Algorithm}

\textbf{(a)} A stable partition compares each nonpivot element with the
pivot once, costing $n-1$ comparisons. Let $s_n$ be the
cost on an already sorted input. Stability leaves both parts sorted, so
\[
s_0=s_1=0,\qquad s_n=2s_{n-1}+n-1=2^n-n-1\quad(n\geq 1).
\]
The first pivot has each rank with probability $1/n$. Conditional on its
rank, stability preserves uniformly random orders within the two parts.
The first recursive pass sorts them; the second therefore costs
$s_k+s_{n-1-k}$. Hence, for $n\geq2$,
\begin{align*}
t_n&=n-1+\frac1n\sum_{k=0}^{n-1}
       (t_k+t_{n-1-k}+s_k+s_{n-1-k})\\
&=n-1+\frac2n\sum_{k=0}^{n-1}t_k
 +\frac2n\left(2^n-1-\frac{n(n+1)}2\right)\\
&=-2+\frac2n\left(2^n-1+\sum_{k=0}^{n-1}t_k\right),
\end{align*}
with $t_0=t_1=0$.
This is exactly the required recurrence.

\textbf{(b)} Multiplying the recurrence by $n$ gives
\[
nt_n=-2n+2(2^n-1)+2\sum_{k=0}^{n-1}t_k
\qquad(n\geq2).
\]
Multiply by $z^{n-1}$ and sum over $n\geq2$. Since $t_0=t_1=0$,
the left-hand side is $T'(z)$. For the prefix-sum term, interchange
the sums to obtain
\[
\sum_{n\geq2}z^{n-1}\sum_{k=0}^{n-1}t_k
=\sum_{k\geq2}t_k\sum_{n\geq k+1}z^{n-1}
=\sum_{k\geq2}\frac{t_kz^k}{1-z}
=\frac{T(z)}{1-z}.
\]
Differentiating the geometric series gives
\[
\sum_{n\geq1}nz^{n-1}=\frac1{(1-z)^2}.
\]
Removing the $n=1$ terms, we get
\begin{align*}
\sum_{n\geq2}nz^{n-1}
&=\frac1{(1-z)^2}-1,\\
\sum_{n\geq2}(2^n-1)z^{n-1}
&=\frac{4z}{1-2z}-\frac{z}{1-z}
=\frac2{1-2z}-\frac1{1-z}-1.
\end{align*}
Substituting these identities and combining fractions yields
\begin{align*}
T'(z)&=\frac{2T(z)}{1-z}
 -2\left(\frac1{(1-z)^2}-1\right)
 +2\left(\frac2{1-2z}-\frac1{1-z}-1\right)\\
&=\frac{2}{1-z}T(z)+\frac{2z}{(1-2z)(1-z)^2}.
\end{align*}

\textbf{(c)} The integrating factor is $(1-z)^2$. Thus
\[
\bigl((1-z)^2T(z)\bigr)'
=\frac{2z}{1-2z}=\frac1{1-2z}-1.
\]
Integrating and using $T(0)=0$ gives
\[
\boxed{T(z)=-\frac{\ln(1-2z)+2z}{2(1-z)^2}.}
\]

\textbf{(d)} Using $-\ln(1-2z)=\sum_{k\geq1}2^kz^k/k$ and
$(1-z)^{-2}=\sum_{j\geq0}(j+1)z^j$, extract coefficients in (c):
\begin{align*}
t_n
&=\frac12\sum_{k=1}^n\frac{(n-k+1)2^k}{k}-n\\
&=\frac{n+1}{2}\sum_{k=1}^n\frac{2^k}{k}
 -\frac12\sum_{k=1}^n2^k-n\\
&=\boxed{\frac{n+1}{2}\sum_{k=1}^n\frac{2^k}{k}-2^n-n+1.}
\end{align*}

\textbf{(e)} Set $j=n-k$ to obtain
\[
\sum_{k=1}^n\frac{2^k}{k}
=\frac{2^n}{n}\sum_{j=0}^{n-1}\frac{2^{-j}}{1-j/n}
=\frac{2^n}{n}\left(\sum_{r=0}^7\frac{M_r}{n^r}+O(n^{-8})\right),
\qquad M_r=\sum_{j\geq0}j^r2^{-j}.
\]
For $j\leq n/2$, expand $(1-j/n)^{-1}$ through degree seven;
the remainder is at most $2(j/n)^8$, whose weighted sum is $O(n^{-8})$.
The tail $j>n/2$ and the tails of the moment sums are exponentially small
times polynomial factors, so they are also $O(n^{-8})$.
Repeatedly applying $x\,d/dx$ to $(1-x)^{-1}$ at $x=1/2$ gives
\[
(M_0,\ldots,M_7)=(2,2,6,26,150,1082,9366,94586).
\]
Substitute into (d). The constant exponential term cancels $2^n$;
the coefficient of $n^{-r}$ inside the parentheses below is
$(M_r+M_{r+1})/4$ for $0\leq r\leq6$. Therefore
\[
\boxed{t_n=\frac{2^{n+1}}n\left(1+\frac2n+\frac8{n^2}
 +\frac{44}{n^3}+\frac{308}{n^4}+\frac{2612}{n^5}
 +\frac{25988}{n^6}+O(n^{-7})\right)-n+1.}
\]

\newpage
\section*{Task 2. Mutual Recursion}

\textbf{(a)} Let $p_n,q_n$ denote the returned values. Both sequences are
nondecreasing, and, for $n\geq3$, each value is at least twice the smaller
value at index $n-3$. Also $p_n,q_n\leq2^n$ by induction. Therefore
\[
2^{\lfloor n/3\rfloor}\leq p_n,q_n\leq2^n,
\]
so both have $\Theta(n)$ bits. Adding the returned integers costs
$\Theta(n)$ bit operations; index arithmetic costs $O(\log n)$.
This justifies the simplified local costs $c_pn,c_qn$.

\textbf{(b)} Set $P(z)=\sum_{n\geq0}T_p(n)z^n$,
$Q(z)=\sum_{n\geq0}T_q(n)z^n$, and
\[
R(z)=\sum_{n\geq3}nz^n=\frac{z^3(3-2z)}{(1-z)^2}.
\]
Summing the recurrences, including the three initial values $1,1,1$, gives
\begin{align*}
P-1-z-z^2&=z(P-1-z)+z^3Q+R,\\
Q-1-z-z^2&=z(P-1-z)+z(Q-1-z)+R.
\end{align*}
Rearranging,
\[
(1-z)P-z^3Q=1+R,\qquad
-zP+(1-z)Q=1-z-z^2+R.
\]
The determinant is $(1-z)^2-z^4=(1-z-z^2)(1-z+z^2)$.
With $D=1-z-z^2$ and $E=1-z+z^2$, solving gives
\[
P=\frac{(1-z)(1+R)+z^3(D+R)}{DE},\qquad
Q=\frac{z(1+R)+(1-z)(D+R)}{DE}.
\]
Substituting $R$ and decomposing into partial fractions yields
\begin{align*}
P(z)&=2+3z+\frac{5+z}{2D}-\frac{3+z}{2E}
      -\frac{z}{(1-z)^2}-\frac2{1-z},\\
Q(z)&=\frac{6+5z}{2D}+\frac{4-3z}{2E}
      -\frac{z}{(1-z)^2}-\frac4{1-z}.
\end{align*}
Use the coefficient identities
\[
\sum_{n\geq0}L_nz^n=\frac{2-z}{D},\quad
\sum_{n\geq0}F_nz^n=\frac zD,\quad
\sum_{n\geq0}\cos(n\pi/3)z^n=\frac{1-z/2}{E},\quad
\sum_{n\geq0}\sin(n\pi/3)z^n=\frac{\sqrt3\,z}{2E}.
\]
For $n\geq2$, coefficient extraction gives
\begin{align*}
T_p(n)&=\frac54L_n+\frac74F_n
 -\frac5{2\sqrt3}\sin\frac{n\pi}3-\frac32\cos\frac{n\pi}3-n-2,\\
T_q(n)&=\frac32L_n+4F_n
 -\frac1{\sqrt3}\sin\frac{n\pi}3+2\cos\frac{n\pi}3-n-4.
\end{align*}
The $T_q$ formula also holds at $n=1$. The stated $T_p$ formula has a
boundary error: it gives $-2$ at $n=1$, whereas $T_p(1)=1$; the $3z$
term supplies the missing $3$.

\textbf{(c)} Memoize each pair (function name, input) while keeping the
original recursive calls. There are $\Theta(n)$ computed states, and the
addition at index $k$ costs $\Theta(k)$ bits of work. Thus the time and
cached-value space are both $\sum_{k=1}^n\Theta(k)=\Theta(n^2)$;
the recursion stack fits within this space bound.

For $\Theta(n)$ extra bits, evaluate the same recurrences bottom-up.
Initialize $p_0=p_1=p_2=q_0=q_1=q_2=1$ and, for $k=3,\ldots,n$, compute
\[
p_k=p_{k-1}+q_{k-3},\qquad q_k=p_{k-1}+q_{k-1}.
\]
Keep only the previous $p$ and the last three $q$'s, computing both new
values before overwriting them. This uses $\Theta(n)$ bits and still
takes $\Theta(n^2)$ time. For $n<3$, simply return $(1,1)$.

\newpage
\section*{Task 3. Akra--Bazzi Recurrences}

Assume positive, bounded base-case costs. For real $x$,
$0\leq x-2\lfloor x/2\rfloor<2$, so $f(x)\in\{0,1\}$.
We bound each recurrence by recurrences with smooth nonnegative tolls;
comparison follows by induction. This avoids applying Akra--Bazzi
directly to the discontinuous function $f$.

For a comparison recurrence $U(n)=\sum_i a_iU(b_in)+g(n)$,
Akra--Bazzi gives
\[
U(n)=\Theta\!\left(n^p\left(1+\int_1^n
\frac{g(u)}{u^{p+1}}\,du\right)\right),
\qquad \sum_i a_ib_i^p=1.
\]
We apply this formula in dependency order.

\textbf{$T_A$.} The characteristic equation
$\frac43 2^{-p}+3^{-p}=1$ has root $p=1$.
The toll lies between $n$ and $n+4n^3$; hence
\[
T_A(n)=\Omega(n\log n),\qquad T_A(n)=O(n^3).
\]
The corresponding integrals have orders $\log n$ and $n^2$.

\textbf{$T_B$.} Let $\beta$ solve $3\,2^{-\beta}+\frac94\,4^{-\beta}=1$.
Then
\[
\beta=\log_2\!\left(\frac{3(1+\sqrt2)}2\right),\qquad
\frac32<\beta<2.
\]
The toll lies between $n^{3/2}$ and $n^{3/2}+3n^3$.
The lower integral is bounded, and the upper is $\Theta(n^{3-\beta})$, so
\[
T_B(n)=\Omega(n^\beta),\qquad T_B(n)=O(n^3).
\]

\textbf{$T_C$.} Here $2\,2^{-p}=1$, so $p=1$.
Using the bounds on $T_A$, the toll is
$\Omega(n\log n)$ and $O(n^3)$.
The comparison integrals are $\Theta(\log^2 n)$ and $\Theta(n^2)$, giving
\[
T_C(n)=\Omega(n\log^2 n),\qquad T_C(n)=O(n^3).
\]

\textbf{$T_D$.} Now $4\,3^{-p}=1$, so $p=\delta=\log_3 4<4/3$.
Since $T_B(2\sqrt n/5)=O(n^{3/2})$, the toll is bounded below by
$n^{4/3}$ and above by $O(n^{3/2})$.
The integrals have orders $n^{4/3-\delta}$ and $n^{3/2-\delta}$, so
\[
T_D(n)=\Omega(n^{4/3}),\qquad T_D(n)=O(n^{3/2}).
\]

\textbf{$T_E$.} Its characteristic equation is
$\frac53 4^{-p}+\frac74 3^{-p}=1$, whose root is $p=1$.
The term $3T_D(n/4)$ gives a lower toll $\Omega(n^{4/3})$;
both $T_C(\sqrt n/2)$ and $T_D(n/4)$ are $O(n^{3/2})$.
Thus the lower and upper integrals have orders $n^{1/3}$ and $n^{1/2}$,
respectively. Therefore the requested bounds are
\[
\boxed{T_E(n)=\Omega(n^{4/3}),\qquad T_E(n)=O(n^{3/2}).}
\]


\end{document}

